\section{A two-prime upper bound for polynomial ratios}
\label{sec:triangle}
Let $r(n)$ be the largest size of a set in $\{0,\ldots,n-1\}$ free of nonconstant three-term arithmetic progressions, with $r(0)=0$. Write $z=1/q$. The one-prime bound of Asada et al.\ is
\[
 U_q=(q-1)\sum_{n\ge1}(r(n)-r(n-1))q^{-n}.
\]
The sum runs over the first lengths at which the maximum size increases. A second linear prime forces strict losses over this bound.

\begin{theorem}
For every prime power $q$, the upper density of any subset of $\mathbb F_q[x]$ avoiding nonconstant polynomial-ratio geometric progressions is at most
\[
 U_q-(1-z)^2(z^8+z^{19}+z^{29}).
\]
In particular the supremum for $q=2$ is at most $907739729/1073741824<0.845398501494$. This answers the smooth-polynomial upper-bound question (Question 6.2).
\end{theorem}

\begin{proof}
Choose distinct monic linear primes $P,Q$. Up to degree $N$, the nonzero polynomials of the form $bP^iQ^j$, for fixed $b$ coprime to $PQ$ with its leading coefficient retained, have indices in
\[
 T_L=\{(i,j)\in\mathbb Z_{\ge0}^2:i+j<L\},\qquad L=N-\deg b+1.
\]
Each row and each column is progression-free in its varying index, since a nonconstant progression on an axis gives a polynomial ratio. Summing just the row bounds permits $R_L=\sum_{k=1}^L r(k)$ positions. The same bound is obtained by columns. Thus the total deficiency $d=R_L-\#\{\text{selected positions}\}$ can be computed either by summing row deficiencies or by summing column deficiencies.

Suppose at least $h$ deficiencies are forced in $T_K$. For $L\ge K$, the last $K$ rows form a translated $T_K$, and the last $K$ columns do also. Their lengths are unchanged by translation. Consequently $d\ge h$. If $d=h$, every row and every column indexed below $L-K$ attains its maximum. For $K=0$ take $h=0$. The following finite interval facts, with indices starting at zero, give three contradictions to equality.
\[
\begin{array}{c|c|c|c|c}
 L&K&i&\text{row }i:\ \text{all maxima include }0&
         \text{column }0:\ \text{all maxima omit }i\\ \hline
 9&0&4 &r(5)=4&r(9)=5\\
 20&9&7&r(13)=7&r(20)=9\\
 30&20&4&r(26)=11&r(30)=12
\end{array}
\]
In each row of the table $i<L-K$, so the two extremizers would disagree at $(i,0)$. Starting with $T_0$ empty and applying the argument successively proves that the deficiency is at least $1$, $2$, $3$ at the displayed lengths. Embedding these three tips in longer triangles therefore gives a penalty of $\mathbf 1_{L\ge9}+\mathbf 1_{L\ge20}+\mathbf 1_{L\ge30}$ in each fiber.

Here is a reproducible finite proof of the interval entries. Generate increasing AP-free subsets by appending a permitted index $t$ and removing from the future choices each $2t-a$ with $a$ already selected. This visits every AP-free subset exactly once. Group by maximum and size, then take unions and intersections of the maximum-cardinality masks across all eligible maxima. The included \texttt{triangle\_check.py} verifies the table with integer arithmetic. For the first row one can check directly that the sole size-four set on five positions is $\{0,1,3,4\}$ and that the size-five sets on nine positions are
\[
 \{0,1,3,7,8\},\quad \{0,1,5,6,8\},\quad
 \{0,1,5,7,8\},\quad \{0,2,3,7,8\}.
\]

To sum the row bounds, merge each row into a one-prime fiber, with $P$ as the varying prime and $P$-free base $bQ^j$. Dividing by $q^{N+1}$, the number of such bases of degree $N-k$ tends to $(1-z)^2 z^k$ at fixed $k$. This follows by counting degree-$N-k$ polynomials and excluding the multiples of $P$. Thus their contribution is bounded in the limit by
\[
 (1-z)^2\sum_{k\ge0}r(k+1)z^k=U_q .
\]
The equality follows by summing differences, and domination follows from $r(k+1)\le k+1$. Each penalty for threshold $L\ge s$ is counted once per two-prime-free $b$ of degree at most $N+1-s$. By inclusion--exclusion its normalized count tends to $(1-z)^2z^{s-1}$. This proves the stated upper bound. For the exact numerical estimate the checker computes $r(n)$ through $n=32$, then bounds the remaining increments by one, giving a tail at $q=2$ of at most $2^{-32}$.
\end{proof}

The degree construction of Asada et al. gives the lower bound $m_q=(1-q^{-2})\prod_{j\ge1}(1+q^{-3^j})$. At $q=2$, retaining just the first four factors shows $m_2>0.845397955517$. The remaining interval for Conjecture 6.3 has width less than $5.46\cdot10^{-7}$. The local calculation uses only ratios that are powers of one of two distinct smooth primes.
